Un skieur chaussé de ses skis (de masse $m=\qty{90}{\kg}$) remonte une piste en
téléski. Le système mécanique étudié est l'ensemble \{skieur~; skis\}. Il y a
des frottements sur la neige. La force de frottement $\vec{f}$ a la même
direction que le vecteur vitesse et le sens inverse~; sa valeur constante vaut
$f=\qty{30}{\N}$.

La remontée en téléski correspond à un déplacement $AB$ de longueur
$\qty{1250}{\m}$~; le plan de remontée forme un angle $\alpha=\ang{20}$ avec
l'horizontale.
\begin{figure}[H]
    \centering
    \includegraphics[width=.9\textwidth]{2025_12_01_2444c970f150c248fc83g-3}
\end{figure}

\begin{enumerate}
    \item Reproduire le schéma et représenter les forces s'exerçant sur le
          système.
    \item Calculer le travail du poids $\vec{P}$ du système lors du déplacement
          considéré. On prendra~: $g=\qty{10}{\N\per\kg}$.
    \item Calculer le travail de la force de frottement $\vec{f}$ lors du
          déplacement considéré.
    \item Dans ces conditions, quel est le travail de la réaction du support
          neigeux sur le système~?
    \item Calculer le travail de la tension du câble qui tire le skieur sachant
          que sa valeur est 675 N et que la direction du câble fait avec le plan
          incliné un angle $\beta=\ang{60}$.
\end{enumerate}

